
Prior work has established that deep learning models lose 10-14\% accuracy when evaluated on ImageNet-V2, an independently-collected variant of ImageNet \citep{recht2019imagenet}. This paper examines whether relative model rankings remain stable despite these accuracy drops. Across 96 models with results on both benchmarks, Kendall-$\tau$ = 0.9647, indicating that the model selected based on ImageNet performance is almost certainly the same model that would be selected based on ImageNet-V2 performance. This finding has practical implications for benchmark-based model selection under distribution shift.

\subsubsection{The Problem}

Research attention is concentrated on a small set of benchmark datasets. The top 10\% of datasets account for 90\% of benchmark usage in machine learning research \citep{koch2021reduced}. For computer vision, ImageNet dominates: a decade of model development has been optimized against its validation set. When Recht et al. (2019) created ImageNet-V2 by replicating the original data collection methodology, models universally experienced 11-14\% accuracy drops, raising concerns about benchmark overfitting.

However, accuracy degradation is not the same as ranking instability. A model ranked 10th on ImageNet might remain ranked 10th on ImageNet-V2, even if both accuracy values are lower. The critical question for practitioners is whether relative model quality transfers, not whether absolute accuracy transfers. If rankings are preserved, then benchmark-based model selection remains valid despite distribution shift.

This distinction has not been systematically examined. Prior work focused on absolute accuracy metrics, documenting drops but not analyzing whether those drops preserve or disrupt the relative ordering of models.

\subsubsection{Approach}

This paper conducts a systematic measurement of ranking stability between ImageNet and ImageNet-V2:

\begin{enumerate}
\item \textbf{Data Collection.} Accuracy results for 96 models evaluated on both ImageNet and ImageNet-V2 are collected from Papers With Code leaderboards.
\end{enumerate}

\begin{enumerate}
\item \textbf{Ranking Correlation.} Kendall-$\tau$ (tau-b) is computed between the two ranking lists. A threshold of $\tau$ < 0.90 is pre-registered as indicating significant ranking shift.
\end{enumerate}

\begin{enumerate}
\item \textbf{Statistical Inference.} 95\% confidence intervals are estimated via bootstrap resampling (10,000 iterations) using the percentile method.
\end{enumerate}

\subsubsection{Key Finding}

The central finding is that accuracy degradation is approximately uniform across models. When all models lose roughly the same percentage of accuracy, their relative rankings are preserved. This explains why $\tau$ = 0.9647: distribution shift affects model accuracy but not model ordering.

\subsubsection{Contributions}

\begin{itemize}
\item \textbf{First systematic ranking stability analysis.} Kendall-$\tau$ between ImageNet and ImageNet-V2 leaderboards is computed across 96 models, yielding $\tau$ = 0.9647 with 95\% CI [0.9454, 0.9795].
\end{itemize}

\begin{itemize}
\item \textbf{Evidence of uniform degradation.} Accuracy drops are approximately uniform (mean 11.68\%, SD 1.87\%), explaining why rankings are preserved despite accuracy loss.
\end{itemize}

\begin{itemize}
\item \textbf{Practical validation of benchmark-based selection.} The results indicate that model selection decisions based on ImageNet generalize to ImageNet-V2-like distributions.
\end{itemize}

